\section{Complexity Analysis}
\label{Section: Convergence analysis}
We now analyze the complexity of the proposed algorithms in terms
of the number of stochastic first-order oracle (SFO) calls
required to find an $\varepsilon$-KKT point of $\calL$.
Given $\varepsilon>0$, a random point $(\bx,\bB,\bLmbd)$
is called an \emph{$\varepsilon$-KKT point} if
\begin{equation}
\label{equ:esp}
\mathbb{E}\big[\operatorname{dist}^{2}\big(\bm{0},\mathcal{K}(\bx,\bB,\bLmbd)\big)\big]
\leq\varepsilon^2.
\end{equation}
\blue{Let $T$ denote the maximum number of iterations. Under the mini-batch sampling scheme
introduced in Section~\ref{Sec: DualHash framework}, let $(\Omega,\mathcal F,\mathbb P)$
be the underlying probability space. For $k\geq1$, define
$\mathcal F_k:=\sigma(\mathcal J_1,\ldots,\mathcal J_{k-1})$ with $\mathcal F_1$ trivial, as
the filtration generated by the past mini-batches. We write
$\mathbb E_k[\cdot]:=\mathbb E[\cdot\mid\mathcal F_k]$ for the conditional expectation
given the past and $\mathbb E[\cdot]$ for the total expectation. All equalities and
inequalities involving $\mathbb E_k[\cdot]$ are understood to hold almost surely.} 

For the complexity analysis, we retain Assumption~\ref{Ass: basic regularity} from
Section~\ref{Section: Problem Reformulation}. \blue{Throughout this section, assume that $h$ is the $\ell_1$ W-regularizer or, more generally, a W-type regularizer with a compact dual domain (\Cref{Section: Problem Reformulation}).} \blue{Let $L_F:=2L_f+\lambda\sqrt{nd}\,L_{\nabla\mathcal{D}}+\frac{\gamma}{n}\bigl(2L_{\nabla\mathcal{D}}C_{\mathcal{D}}+2G_{\mathcal{D}}^2+1\bigr)$. Under the prescribed initialization and step sizes, $F$ is $L_F$-smooth along the iterates (\blue{\Cref{app:smoothness-discussion}}), which is essential for our analysis.} We then impose the following stochastic-oracle assumption on Problem \eqref{Prop: block-wise regularized problem} used in this section.
\begin{assumption}
\label{Ass: main assumptions}
\leavevmode
There exists $\sigma^2\geq0$ such that, for all $\bx\in\RR^{q}$ and $\bB\in\RR^{n\times d}$,
    \begin{align}
        \label{Eq: finite-sum sampling identities}
        &\blue{\frac{1}{n}\sum_{i=1}^{n}\nabla_{\bx} F_i(\bx, \bb_i)
        = \nabla_{\bx} F(\bx, \bB)},\\
        &\blue{\frac{1}{n}\sum_{i=1}^{n}
        \left\|\nabla_{\bx} F_i(\bx, \bb_i)
        - \nabla_{\bx} F(\bx, \bB)\right\|^2
        \leq \sigma^2}.
    \end{align}
\end{assumption}

Our analysis aims to bound $\mathbb{E}\big[\operatorname{dist}^{2}\big(\bm{0},\mathcal{K}(\bx,\bB,\bLmbd)\big)\big]$. However,
\blue{the dual update can increase $\calL$, so $\calL$ need not decrease along the iterates and standard descent arguments do not apply directly}~\cite{xu2023momentum,zhang2022stochastic}.
To this end, we construct suitable Lyapunov functions (see \eqref{Eq: the Lyapunov function in SGDM} and \eqref{Eq: the Lyapunov function in STORM}) to obtain the approximate descent property. We first analyze the dual dynamics by establishing bounds on $\|\bLmbd^{k+1}-\bLmbd^k\|^2$ in terms of the primal residuals, enabling the Lyapunov construction \blue{(\Cref{Lemma: dual boundedness})}.

Next, we analyze how the Lagrangian function $\calL$ changes
during the updates of primal and dual variables in one iteration. The total change can be decomposed as:
\begin{align}
    \label{Eq: One-iteration progress decomposition} & \calL(\bx^{k+1}, \bB^{k+1}, \bLmbd^{k+1}) - \calL(\bx^{k}, \bB^{k}, \bLmbd^{k}) \notag                                                                                                                                                                                       \\
    =                                                & \,\underbrace{\calL(\bx^{k+1}, \bB^{k+1}, \bLmbd^{k+1}) - \calL(\bx^{k+1}, \bB^k, \bLmbd^k)}_{\text{one-iteration progress on $\bB$ and $\bLmbd$}}+ \underbrace{\calL(\bx^{k+1}, \bB^{k}, \bLmbd^{k}) - \calL(\bx^{k}, \bB^k, \bLmbd^k)}_{\text{one-iteration progress on $\bx$}}.
\end{align}
Since the updates of $\bB$ and $\bLmbd$ are common to Algorithms~\ref{alg:our_algorithm_sgdm} and~\ref{alg:our_algorithm_storm}, \blue{\Cref{Lemma: One-iteration progress of B and Lambda}} bounds their one-step change. For notational simplicity, we define
\[
    \Delta_{\bx}^{k}= \frac{\|\bx^{k}- \bx^{k-1}\|^{2}}{2},\quad \Delta_{\bm{B}}^{k}
    = \frac{\|\bm{B}^{k}- \bm{B}^{k-1}\|^{2}}{2}.
\]

We now analyze the stochastic $\bx$-updates, which \blue{introduce} algorithm-specific challenges:
as \eqref{Eq: dual boundedness} and \blue{\eqref{item: B and Lambda}} show, randomness in the $\bx$-updates propagates to the $\bB$- and $\bLmbd$-updates, necessitating careful parameter tuning and the construction of Lyapunov functions. We present the key lemmas and complexity results of both algorithms below.
\subsection{Oracle Complexity of DualHash-StoM}
\label{Section: complexity of StoM}
\blue{We now present the oracle complexity of DualHash-StoM for finding an $\varepsilon$-KKT point, with the detailed analysis deferred to \Cref{app:stom-complexity}.} To handle the oscillation of \blue{$\calL$}, we \blue{first set the parameters of \Cref{alg:our_algorithm_sgdm} as}
\begin{equation}
    \label{Eq: the parameters conditions of dualhash-sgdm in suppl}
    \eta_{k}= \frac{\eta}{L_{F}},\quad \alpha_{k}= \alpha,\quad \beta_{k}= \beta,\quad \blue{0<\tau<\frac{\sqrt{30}-2}{13L_F}},\quad \blue{\eta=\frac{(1 - 2\alpha - \nu)/2}{\frac{\bar{L}_{F}}{L_{F}} + \frac{(1 + \nu)\beta^{2}}{(1 - \alpha)\alpha}}},
\end{equation}
where \blue{$\bar{L}_F:=2L_F+K_3+K_4$ with $K_1,\ldots,K_4$ given in \Cref{Lemma: One-iteration progress of B and Lambda}, and} \blue{$\nu\in(0,1)$, $\alpha\in(0,\frac{1-\nu}{2})$, $\beta\in(0,1)$ are constants independent of $T$.} We \blue{then} introduce the following Lyapunov function and sequence to establish an in-expectation approximate sufficient descent property for \Cref{alg:our_algorithm_sgdm}:
\begin{align}
    \Psi_m(\mathcal{Q}^{k}) &= \mathcal{L}(\bx^{k}, \bm{B}^{k}, \bLmbd^{k})
        - C_{1}\Delta_{\bB}^{k+1} + C_{2}\Delta_{\bB}^{k}
        - C_{3}\Delta_{\bm{x}}^{k+1} + C_{4}\Delta_{\bm{x}}^{k},
        \label{Eq: the Lyapunov function in SGDM} \\
    \mathcal{Q}^{k} &:= \{\bx^{k}, \bB^{k}, \bm{\Lambda}^{k},
        \bB^{k+1}, \bB^{k-1}, \bx^{k+1}, \bx^{k-1}\}, \notag
\end{align}where the positive coefficients $C_1, C_2, C_3, C_4$ (independent of $k$) are chosen as:
\begin{align}
    \label{Eq: the coefficients of the Lyapunov function in SGDM}
    \begin{cases}
        C_{1} &= K_{1},\\
        C_{2} &= \frac{1}{2}(K_{2} - K_{1} + K_{3}),\\
        C_{3} &= K_{4},\\
        C_{4} &= \frac{\bar{L}_F + \frac{2L_F \beta^2}{\alpha}}{1 - 2\alpha - \nu}.
    \end{cases}
\end{align}\blue{$\Psi_m$ is bounded below by some $\Psi_m^*>-\infty$ independent of $T$ (see \Cref{Lemma: lower bound of Psi_m}); for a deterministic initial point, define $\Delta_1:=\calL(\bx^1,\bB^1,\bLmbd^1)-\Psi_m^*$.}
The construction yields the following in-expectation approximate sufficient
descent property of \blue{$\{\Psi_m(\mathcal{Q}^{k})\}$}.
\begin{lemma}
    \label{Lemma: approximate sufficient descent property in SGDM} Under Assumptions~\ref{Ass: basic regularity} and~\ref{Ass: main assumptions} and the parameter conditions \eqref{Eq: the parameters conditions of dualhash-sgdm in suppl}, \blue{for all $k \geq 1$},
    \begin{align}
        \label{Eq: the descent property of Lyapunov func in SGDM}\EE [\Psi_m(\mathcal{Q}^{k+1})] - \EE[\Psi_m(\mathcal{Q}^{k})] \leq -C_{\bB}\EE[\Delta_{\bB}^{k+1}+ \Delta_{\bB}^{k}] -C_{\bx}\EE [\Delta_{\bm{x}}^{k+1}+ \Delta_{\bm{x}}^{k}] + \frac{\sigma^{2}}{b_{k}L_{F}},
    \end{align}
    where the descent coefficients are given by
    \begin{align}
        \label{Eq: the descent coefficients of potential function in SGDM}\begin{cases}C_{\bB}&= C_{2}- K_{3}= \frac{K_2 - K_3 -K_1}{2}>0, \\ C_{\bx}&= \nu C_{4}>0.\end{cases}
    \end{align}
\end{lemma}
\blue{The proof is provided in \Cref{app:stom-descent-proof}}, and the positivity of the coefficients in
\eqref{Eq: the coefficients of the Lyapunov function in SGDM} and the descent
constants in \eqref{Eq: the descent coefficients of potential function in SGDM}
\blue{are} established in \Cref{lem:lyapunov-coef-positivity}.

We now present the convergence rate of \Cref{alg:our_algorithm_sgdm} under the parameter conditions \eqref{Eq: the parameters conditions of dualhash-sgdm in suppl}, which guides the selection of the batch size $b_k$.
\begin{lemma}
\label{lemma: convergence rate sgdm} Under the conditions of \Cref{Lemma:
    approximate sufficient descent property in SGDM}, for any $T \geq 1$,
    \blue{let $R$ be drawn uniformly from $\{1,\ldots,T\}$, independently of the mini-batches. Then} $(\bx^{R}, \bB
    ^{R}, \bLmbd^{R})$ satisfies \eqref{equ:esp} with
    \begin{align}
        \label{Eq: the convergence rate of sgdm}\mathbb{E}[\mathrm{dist}^{2}(\blue{\bm{0}, \mathcal K}(\bx^{R}, \bB^{R}, \bLmbd^{R}))] \leq \frac{L_{\Delta}\Delta_{1}}{T}+ \blue{\frac{C_{\sigma}\sigma^{2}}{T}\sum_{k=1}^{T}\frac{1}{b_{k}}},
    \end{align}\blue{where $L_{\Delta}:=\frac{\max\{S_1,S_2,S_3,S_4\}}{\min\{C_{\bB},C_{\bx}\}}$ and $C_{\sigma}:=\frac{L_{\Delta}}{L_F}+3$ are constants independent of $T$, with $S_1,\ldots,S_4$ defined in \eqref{Eq: the coefficients of stationarity error bound in SGDM}.}
\end{lemma}
\blue{To achieve the desired oracle complexity, we choose the batch size according to \eqref{Eq: the convergence rate of sgdm} as}
\begin{align}
\label{Eq: the parameters conditions of dualhash-sgdm}
    \blue{b_k=b=\lceil c_bT\rceil\in\{1,\ldots,n\},}
\end{align}
\blue{where $c_b>0$ is a constant independent of $T$, and $1\le\lceil c_bT\rceil\le n$.} The complexity result is established in \Cref{Theorem: the oracle complexity in SGDM}, with the proof in \Cref{app:stom-oracle-proof}.
\begin{theorem}
% [Oracle complexity of \Cref{alg:our_algorithm_sgdm}]
    \label{Theorem: the oracle complexity in SGDM}
Under \blue{Assumptions~\ref{Ass: basic regularity} and}~\ref{Ass: main assumptions}, let \(\{\bx^k, \bB^k, \bLmbd^k\}_{k \in \NN}\) be the iterate sequence from \Cref{alg:our_algorithm_sgdm} with the parameters satisfying \blue{\eqref{Eq: the parameters conditions of dualhash-sgdm in suppl} and} \eqref{Eq: the parameters conditions of dualhash-sgdm}. Then, \blue{for $R$ drawn uniformly from $\{1,\dots,T\}$ independently of the mini-batches,} \( (\bx^R, \bB^R, \bLmbd^R )\) satisfies \eqref{equ:esp} \blue{with}
\begin{align}
    \label{Eq: the iterate number upper bound of algorithm SGDM}
        T \blue{\geq} \left\lceil \frac{L_{\Delta} \Delta_1+ C_{\sigma}\frac{\sigma^2}{c_b}}{\varepsilon^2} \right\rceil,
\end{align}
\blue{where $L_{\Delta}$ and $C_{\sigma}$ are given in \Cref{lemma: convergence rate sgdm}.} \blue{Consequently}, the total number of SFO calls is of order \(\calO(\varepsilon^{-4})\).
\end{theorem}
\subsection{Oracle Complexity of DualHash-StoRM}
\blue{We now present the oracle complexity of DualHash-StoRM for finding an $\varepsilon$-KKT point, with the detailed analysis deferred to \Cref{app:storm-complexity}.} Before proceeding, we introduce an additional assumption, also known as the mean-squared smoothness
assumption, widely used in variance reduction methods~\cite{arjevani2023lower,fang2018spider,tran2020hybridvrsgd,driggs2021stochastic}.
\begin{assumption}
\label{Ass: the mean squared smoothness of F}
There exists $L_{F} > 0$ such that \blue{for all $\bx, \bx' \in \RR^{q}$ and $\bB, \bB' \in \RR^{n\times d}$}, it holds that
\begin{align}
    \label{eq:mss}
    \blue{\frac{1}{n}\sum_{i=1}^{n}}
    \left\| \nabla_{\bx} F_{i}(\bx,\bb_i) - \nabla_{\bx} F_{i}(\bx',\bb_i')\right\|^{2}
    \leq L^{2}_F \left[ \|\bx - \bx'\|^{2} + \|\bB - \bB'\|^{2} \right].
\end{align}
\end{assumption}\blue{By Jensen's inequality, \eqref{eq:mss} implies that $\nabla_{\bx}F$ is $L_F$-Lipschitz in $(\bx,\bB)$. The converse need not hold.}

\blue{With a slight abuse of notation, $\be^k$ denotes the algorithm-specific gradient error, and $\mathcal{Q}^k$ the corresponding iterate tuple.} Let us denote the gradient error \[\be^{k} = \blue{\bg_s^{k}} - \nabla_{\bx}F(\bx^{k}, \bB^{k})
,\]where $\blue{\bg_s^k}$ is defined in \eqref{Eq: the stochastic approximate in STORM}. \blue{We first set the parameters of \Cref{alg:our_algorithm_storm} as}
\begin{align}
    &\eta_k \blue{\equiv\eta_T}:= \frac{\eta}{\tilde{L}_{F} T^{1/3}},\quad
    \rho_k \blue{\equiv\rho_T}:= \frac{8\rho\eta^2}{T^{2/3}},\quad
    b_k = \begin{cases} b_1, & k = 1, \\ b, & k \geq 2, \end{cases}
    \label{Eq: the parameters conditions in dualhash-storm}\\
    &\blue{0<\tau<\frac{\sqrt{30}-2}{13L_F}},\quad
    \blue{\rho\ge\frac{2L_F^{2}(1+G)}{b\tilde{L}_F^{2}}},\quad
    \blue{0<\eta\le\min\Bigl\{\frac12,\frac{1}{4\sqrt{\rho}}\Bigr\}},
    \label{Eq: the constants for storm}
\end{align} 
where \blue{$\tilde{L}_F:=L_F+K_3+K_4$,} $\eta$ and $\rho$ are positive constants independent of $T$, \blue{$b\in\{1,\ldots,n\}$ is an integer, and $G:=c_GG_{\mathcal{D}}^{2}$ for a universal constant $c_G>0$ specified in \blue{\Cref{lem:auxiliary-tracking}}.} We \blue{then} introduce the following Lyapunov function of \Cref{alg:our_algorithm_storm}:
% and sequence specific to \Cref{alg:our_algorithm_storm}:
\begin{align}
    \Psi_r(\mathcal{Q}^{k}) &= \calL(\bx^{k}, \bm{B}^{k}, \bLmbd^{k})
    - C_{1}\Delta_{\bB}^{k+1}+ C_{2}\Delta_{\bB}^{k}- C_{3}\Delta_{\bm{x}}^{k+1},
    \label{Eq: the Lyapunov function in STORM}
    % \\
    % \{ \calQ^{k}\}_{k \in \NN} &= \{\bx^{k}, \bB^{k}, \bm{\Lambda}^{k},
    % \bB^{k+1}, \bB^{k-1}, \bx^{k+1}\}_{k \in \NN},
    % \label{Eq: the Lyapunov sequence in STORM}
\end{align}
where $C_1$, $C_2$, and $C_3$, which are independent of $k$, are chosen as in
\eqref{Eq: the coefficients of the Lyapunov function in SGDM}. 
\blue{Similarly, \blue{the augmented potential $\EE[\Psi_r(\calQ^k)]+\frac{\eta_T}{\rho_T}\EE\|\be^k\|^2$} is bounded below by some $\Psi_r^{*}>-\infty$ independent of $T$ (see \Cref{Lemma: lower bound of W}); define $\tilde{\Delta}_{1}:=\calL(\bx^{1},\bB^{1},\bLmbd^{1})-\Psi_r^{*}$.}

In the next lemma, we establish an \blue{in-expectation} approximate sufficient descent property
of $\Psi_r$, with the proof provided in \blue{\Cref{app:storm-descent-proof}}.
\begin{lemma}
\label{Lemma: approximate sufficient descent property storm}
    \blue{Under \Cref{Ass: basic regularity}(ii),} with \blue{$0<\tau<\frac{\sqrt{30}-2}{13L_F}$ and $\eta_{k}^{-1}- \tilde{L}_{F}>0$}, for all $k \geq 1$,
    \begin{align}
        \blue{\EE}[\Psi_r(\mathcal{Q}^{k+1})] - \blue{\EE}[\Psi_r(\mathcal{Q}^{k})]
        &\leq -C_{\bB}\blue{\EE}[\Delta_{\bB}^{k+1}+\Delta_{\bB}^{k}]-C_{\bx}^{k}\blue{\EE}[\Delta_{\bx}^{k+1}]+\frac{\eta_k}{2}\blue{\EE}[\|\be^{k}\|^{2}],
        \label{Eq: the descent property of Lyapunov func in storm}
    \end{align}
    where the descent coefficients are given by:
    \begin{align}
        % \label{Eq: the descent coefficients of potential function}
        \begin{cases}C_{\bB}&= C_{2}- K_{3}= \frac{K_2 - K_3 -K_1}{2}>0, \\ C_{\bx}^{k}&=\eta_{k}^{-1}- \tilde{L}_{F} \blue{> 0}.\end{cases} \notag
    \end{align}
\end{lemma}

We now analyze the convergence rate of \Cref{alg:our_algorithm_storm}\blue{, which determines the choice of the initial batch size $b_1$}.
\begin{lemma}
    \label{Lemma: convergence rate of storm} Under \blue{Assumptions~\ref{Ass: basic regularity}, \ref{Ass: main assumptions}, and~\ref{Ass: the mean squared smoothness of F}} and the parameter conditions \eqref{Eq: the parameters conditions in dualhash-storm}--\eqref{Eq: the constants for storm}, for any $T \geq 1$, \blue{let $R$ be drawn uniformly from $\{1,\ldots,T\}$, independently of the mini-batches. Then} $(\bx^{R}, \bB^{R}, \bLmbd^{R})$ satisfies \eqref{equ:esp} with
    \begin{align}
    \label{Eq: the convergence rate of storm}\EE[\mathrm{dist}^{2}(\blue{\bm{0}, \mathcal K}(\bx^{R}, \bB^{R}, \bm{\Lambda}^{R}))] \leq \frac{\tilde{L}_{\Delta}\tilde{\Delta}_{1}}{T^{2/3}}+\blue{\frac{\tilde{C}_{\sigma}\sigma^{2}}{4\rho\eta^{2}b_1T^{1/3}}}+\frac{32\rho\eta^{2}\tilde{C}_{\sigma}\sigma^{2}}{b\,T^{2/3}},
    \end{align}
    where
    \blue{\[
    \tilde{L}_{\Delta}:=\max\Bigl\{\frac{S_1}{C_{\bB}},\frac{S_2}{C_{\bB}},\frac{48\tilde{L}_F}{\eta}+\frac{16\eta L_F^{2}(1+3\tau^{2})}{\tilde{L}_F}\Bigr\},\qquad \tilde{C}_{\sigma}:=\frac{\eta\tilde{L}_{\Delta}}{2\tilde{L}_F}
    \]}
    are constants independent of $T$, with $S_1,S_2$ defined in \eqref{Eq: the coefficients of stationarity error bound in SGDM}.
\end{lemma}

\blue{As noted in~\cite{xu2023momentum}, a moderately large initial batch size improves the complexity order at a negligible cost.} \blue{To achieve the desired oracle complexity, we choose $b_1=\lceil c_bT^{1/3}\rceil$ with $c_b>0$ independent of $T$ and $1\le b_1\le n$.} \blue{The resulting complexity is given in \Cref{Theorem: the oracle complexity in STORM}, whose proof appears in \Cref{app:storm-oracle-proof}.}
\begin{theorem}
\label{Theorem: the oracle complexity in STORM}
Under \blue{Assumptions~\ref{Ass: basic regularity}, \ref{Ass: main assumptions}, and~\ref{Ass: the mean squared smoothness of F}}, let \(\{\bx^k, \bB^k, \bLmbd^k\}_{k \in \NN}\) be the iterate sequence from \Cref{alg:our_algorithm_storm} with the parameters satisfying \eqref{Eq: the parameters conditions in dualhash-storm}--\eqref{Eq: the constants for storm} \blue{and the above choice of $b_1$.} Then, \blue{for $R$ drawn uniformly from $\{1,\dots,T\}$ independently of the mini-batches,} \( (\bx^R, \bB^R, \bLmbd^R )\) satisfies \eqref{equ:esp} whenever
\begin{align}
    \label{Eq: the iterate number upper bound of algorithm VR}
        T \blue{\geq} \left\lceil \frac{\Bigl(\tilde{L}_{\Delta}\tilde{\Delta}_{1}+\tilde{C}_{\sigma}\bigl(\frac{\sigma^{2}}{4\rho\eta^{2}c_b}+\frac{32\rho\eta^{2}\sigma^{2}}{b}\bigr)\Bigr)^{3/2}}{\varepsilon^3} \right\rceil,
\end{align}
with \blue{$\tilde{L}_{\Delta}$ and $\tilde{C}_{\sigma}$ from \Cref{Lemma: convergence rate of storm}}. \blue{Consequently,} the total number of SFO calls is of order \(\calO(\varepsilon^{-3})\).
\end{theorem}
\begin{remark}
\begin{enumerate}[label=(\roman*), leftmargin=2em]
\item The results in Theorems \ref{Theorem: the oracle complexity in SGDM} and \ref{Theorem: the oracle complexity in STORM} \blue{establish stochastic first-order complexity bounds for the proposed multi-block primal-dual scheme, in a setting} where existing work has largely focused on one-block~\cite{xu2023momentum,shi2025momentum} or deterministic \blue{methods}~\cite{ADMMNolinear2024}.
    \item  \Cref{alg:our_algorithm_sgdm} \blue{under the parameter conditions \eqref{Eq: the parameters conditions of dualhash-sgdm in suppl} and} \eqref{Eq: the parameters conditions of dualhash-sgdm} attains both oracle and sample \(\mathcal{O}(\varepsilon^{-4})\) complexities, which matches that of stochastic first-order methods (SFOMs) for one-block stochastic optimization~\cite{Ghadimi2013MinibatchSA}. Meanwhile, \Cref{alg:our_algorithm_storm} with the refined constant parameters from \eqref{Eq: the parameters conditions in dualhash-storm} attains both oracle and sample $\mathcal{O}(\varepsilon^{-3})$ complexities. \blue{The StoRM $\varepsilon^{-3}$ order matches the best-known lower-bound order for smooth stochastic nonconvex optimization~\cite{arjevani2023lower} in the stochastic-oracle regime, for the single-block smooth stationarity measure considered there.}
    \item \blue{Since mini-batches are drawn without replacement, the gradient error is bounded by $\frac{n-b}{n-1}\cdot\frac{\sigma^2}{b}$, which vanishes at $b=n$. Hence, with $b=\min\{\lceil c_bT\rceil,n\}$ and $b_1=\min\{\lceil c_bT^{1/3}\rceil,n\}$, the guarantees of \Cref{Theorem: the oracle complexity in SGDM,Theorem: the oracle complexity in STORM} hold for every $\varepsilon>0$, with $\mathcal{O}(\min\{\varepsilon^{-4},n\varepsilon^{-2}\})$ and $\mathcal{O}(\varepsilon^{-3})$ SFO calls, respectively.}
\end{enumerate}
\end{remark}

\subsection{\blue{Oracle Complexity Guarantee for Problem~\eqref{Eq: the primal deep hashing problem}}}
\label{Sec: guarantees original objective}
\blue{We now present, with \blue{proof in \Cref{app:original-objective}}, the oracle complexity guarantee of DualHash for the original Problem~\eqref{Eq: the primal deep hashing problem}, which follows from combining the KKT-residual bounds in \Cref{Theorem: the oracle complexity in SGDM,Theorem: the oracle complexity in STORM} with \Cref{prop:original-certificate}.}

\begin{corollary}\label{cor:original-stationarity}
Under the assumptions of \Cref{Theorem: the oracle complexity in SGDM} (resp.\ \Cref{Theorem: the oracle complexity in STORM}), \blue{let $(\bx^R,\bB^R,\bLmbd^R)$ be the output of \Cref{alg:our_algorithm_sgdm} (resp.\ \Cref{alg:our_algorithm_storm}), obtained with} $\calO(\varepsilon^{-4})$ (resp.\ $\calO(\varepsilon^{-3})$) SFO calls. Then $\bx^R$ is $(\delta_R,\varepsilon_R)$-near-stationary for \blue{Problem~\eqref{Eq: the primal deep hashing problem}}, where \blue{$\delta_R$ and $\varepsilon_R$ are the bounds of \Cref{prop:original-certificate} at $(\bx^R,\bB^R,\bLmbd^R)$ with $L=L_0$, and satisfy
\begin{align*}
\EE\,\delta_R &\le\sqrt{\frac{2}{L_0}}\Bigl(\sqrt{\kappa}\,\varepsilon+\sqrt{\EE\,h\bigl(\mathcal{D}(\bx^R)\bigr)}\Bigr),\\
\EE\,\varepsilon_R &\le\Bigl(\sqrt{1+G_{\mathcal{D}}^2}+2\sqrt{2\kappa L_0}\Bigr)\varepsilon+2\sqrt{2L_0\,\EE\,h\bigl(\mathcal{D}(\bx^R)\bigr)},
\end{align*}
with $\kappa$ and $L_0$ given in \Cref{lem:smooth-lower-bound} and \Cref{prop:original-certificate}.}
\end{corollary}

\blue{\Cref{cor:original-stationarity} makes \Cref{prop:original-certificate} concrete for the algorithm output: the optimization error is reduced to $\calO(\varepsilon)$ at the oracle complexities of \Cref{Theorem: the oracle complexity in SGDM,Theorem: the oracle complexity in STORM}, whereas the structural error is determined by the quantization error of the output, which is inherent to continuous relaxation (\Cref{Section: Introduction}). In particular, for the $\ell_1$ W-regularizer, $\bx^R$ is $(\calO(\varepsilon),\calO(\varepsilon))$-near-stationary for Problem~\eqref{Eq: the primal deep hashing problem} whenever $\lambda\,\EE\,Q_1(\bx^R)=\calO(\varepsilon^2)$, i.e., when the network outputs are close enough to binary codes, which can be monitored in practice by computing $Q_1(\bx^R)$. It is worth emphasizing that, with our reformulation and algorithm, DualHash attains the smallest quantization error among the compared methods on the challenging multi-label NUS-WIDE dataset (\Cref{Subsection: Further analysis}), so the structural error remains small in practice and the guarantee for the original objective is practically meaningful.}

